\section{Experimental setup and instance designs}
\label{sec:saa}

This section describes the experimental design used to find a suitable sampling
strategy and scenario count that balance outcome coverage with computational
tractability. We first describe how scenarios are constructed from the data,
then the three sampling strategies evaluated, and finally the overall
experimental design including parameter justification.

\subsection{Scenario Generation from Data}
\label{sec:scenario_data}

The dataset contains over $20{,}000$ patient treatment records, each including
specialty $q \in \{\text{CHI},\text{ORT},\text{KNO}\}$, surgery type, OR
identifier, waiting list identifier, and both planned and realised surgery
duration. In total, 186 unique surgery types across 12 ORs are recorded. For a
subset of records, the session identifier and within-session sequence are also
provided.

Two observations from the data motivate our scenario generation approach:
\begin{enumerate}
  \item Realised durations differ significantly across surgery types, making
        surgery type a strong predictor of duration.
  \item Planned durations outperform surgery-type historical averages as point
        estimates: the mean absolute error using planned durations is $12.8\%$
        lower than when using surgery-type averages.
\end{enumerate}

We therefore model the \emph{planned-duration error}
$\varepsilon_p = d_p^{\text{real}} - d_p^{\text{plan}}$ rather than the
realised duration directly, preserving the information content of the planned
duration. For each surgery type $\tau$, we estimate parameters
$\mu_\tau$ and $\sigma_\tau^2$ as the mean and variance of the observed errors
$\varepsilon_p$ over all historical patients of type $\tau$. Realised durations
are then sampled as:
%
\begin{equation}
  d_{ps} \sim
    \mathrm{LogNormal}\!\left(d_p^{\text{plan}} + \mu_\tau,\; \sigma_\tau^2
    \right),
  \qquad \forall\, p \in P,\; s \in S.
  \label{eq:duration_dist}
\end{equation}
%
Sampling the error rather than the realized duration directly ensures that surgery-duration stochasticity
does not absorb known variance that is captured in the planned
duration. To generate a scenario $s$, one planned duration $d_{ps}$ is drawn per
patient $p$ from \eqref{eq:duration_dist}. All patient durations within a
scenario are drawn independently; scenarios are also drawn independently of
each other.

To avoid sampling clinically implausible durations from the unbounded tails of
the lognormal distribution, each draw is truncated to the interval
$[\mu_p - 2\sigma_p,\, \mu_p + 2\sigma_p]$. This
reflects the practical observation that a procedure is extremely unlikely to
take a drastically shorter or longer time than its planned duration suggests,
and that retaining such extreme outliers would distort both the SAA training set
and the out-of-sample evaluation. 

\subsection{Sampling Strategies}
\label{sec:sampling}

We evaluate three sampling strategies to assess which provides the best
solution quality for a given sample size $N$.

\subsubsection{Random Sampling (RS)}
Each duration $d_{ps}$ is drawn independently and uniformly at random from
\eqref{eq:duration_dist}. Scenario probabilities are set to $\pi_s = 1/N$ for
all $s \in S$. RS is unbiased and asymptotically consistent, but with small
$N$ some regions of the distribution may be over- or under-represented, leading
to higher variance across replications.

\subsubsection{Latin Hypercube Sampling (LHS)}
LHS stratifies the marginal cumulative distribution of each patient's duration
into $N$ equally probable intervals $[(k-1)/N,\; k/N]$ for $k = 1, \dots, N$.
One quantile value is drawn uniformly from each interval per patient, and the
$N$ values are then randomly permuted across scenarios to form the scenario
matrix $d_{ps}$. Scenario probabilities are $\pi_s = 1/N$. By construction,
each marginal distribution is evenly covered, yielding lower variance than RS
for the same $N$~\cite{McKay1979}.

\subsubsection{Importance Sampling (IS)}

Importance Sampling (IS) is a variance-reduction technique that draws more
frequently from regions that disproportionately influence the objective and
corrects the bias through likelihood-ratio reweighting. In appointment
scheduling the extreme duration scenarios drive overtime, idling, and waiting costs, yet
random sampling assigns them low probability, requiring many scenarios for
robust solutions for these outlier cases.

To include such tail events we apply the mean-shifting strategy
of~\cite{aditri-2026}: the scenario set $S$ is split into three equal subsets
drawn from lognormals with mean parameters $\mu_\tau - k\sigma_\tau$, $\mu_\tau$,
and $\mu_\tau + k\sigma_\tau$ ($k > 0$). The negative shift targets early
finishes (idling) and the positive shift targets overruns (overtime and
downstream waiting). To keep the SAA objective unbiased, each scenario $s$ is
reweighted by its likelihood ratio under the nominal versus the shifted
distribution it was drawn from:
%
\begin{equation}
  \pi_s
    = \frac{\prod_{i} f(d_{is} \mid \mu_\tau,\, \sigma_\tau)}
           {\prod_{i} f(d_{is} \mid \mu_\tau + \delta_s \sigma_\tau,\, \sigma_\tau)},
  \label{eq:is_weights}
\end{equation}
%
where $f(d_{is} \mid \mu_\tau,\, \sigma_\tau)$ is the nominal lognormal density,
$\delta_s \in \{-k,\, 0,\, k\}$ the shift applied to scenario $s$'s subset, and
the product runs over all surgeries $i$ (the joint density factorises as
durations are independent). Scenarios consistent with the nominal distribution
receive weights near one, while strongly shifted ones receive proportionally
lower weights.

Consequently the tail scenarioscarry small weights, so the effective sample size
$N_{\text{eff}} = 1 / \sum_s \pi_s^2$ is below $N$, and for small $N$ IS may need
more scenarios than RS for a stable estimate. The compensating benefit is
robustness: by guaranteeing that worst-case realisations appear in every sample,
IS yields policies that hedge more conservatively against tail events
(Section~\ref{sec:results}).

\subsection{Experimental Setup}
\label{sec:exp_setup}
In this section, the parameters and used evaluation sample is discussed and the experiments are defined.
\paragraph{closing time}
Based on a 8 hour working day we set the $t^{start} = 480$ and $t^{close} = 960$. Our experiments prove that these settings are an effective balance between fitting all patients in a day and needing overtime in a number of scenarios.
\paragraph{Evaluation sample}
Solution quality is assessed on an out-of-sample evaluation set of size
$N' = 10{,}000$ scenarios generated by RS from the treatment duration distribution
\eqref{eq:duration_dist}. Because the evaluation set is large, RS provides
near-complete coverage of the outcome space and serves as an approximation for the true
expected cost.

\paragraph{SAA parameters.}
We use $M = 10$ independent replications per sample size $N$. The sample size starts from an initial $N = 10$ and increases in steps of $5$. Rather than estimating the true
optimum $v^*$ to bound the optimality gap, we monitor the variance and mean of the
out-of-sample objective across replications through a $95\%$ confidence interval
on the out-of-sample mean cost:
%
\begin{equation}
  \bar{v} \pm z_{0.975}\,\frac{\hat{\sigma}_M}{\sqrt{M}},
  \label{eq:ci}
\end{equation}
%
where $\bar{v}$ and $\hat{\sigma}_M$ are the sample mean and standard
deviation of the $M$ out-of-sample evaluations, and $z_{0.975} = 1.96$.
For each sample size we record this interval and plot it against $N$; the
appropriate $N$ is then selected based on the convergence plot as
the point beyond which the objective stabilises.

\paragraph{Experiments}

\begin{table}[H]
  \centering
  \caption{Experimental settings of the base instance (waiting list~4).}
  \label{tab:instance}
  \begin{tabular}{lr}
    \toprule
    Setting & Value \\
    \midrule
    Patients    & 33 \\
    Sessions    & 6  \\
    Specialties & 3  \\
    Scenarios   & 10 \\
    \bottomrule
  \end{tabular}
\end{table}

Steps~1 and~2 use the waiting-list~4 patient set, whose six documented OR
sessions carry session and position labels that fix the assignment (and, for
Step~1, the sequencing). The characteristics of this waiting list are shown in Table~\ref{tab:instance}. Step~3 uses four patient sets (waiting lists~4, 7, 10,
and~13), each spanning at least six sessions, to assess solution quality across
patient populations. Each sample is solved with all three sampling strategies,
sweeping $N$ from 10 in steps of 5, using the $M = 10$ replications and the
$N' = 1{,}000$ evaluation set described above and the convergence
criterion~\eqref{eq:ci}.
